\section{Булевы функции. Функции одной и двух переменных, функции трех переменных. Существенные и фиктивные переменные.}

\subsection*{Определение булевой функции}

\paragraph{Определение}
Булевой функцией (функцией алгебры логики) от $n$ переменных называется функция $f: E_2^n \to E_2$, где $E_2 = \{0, 1\}$.
\begin{itemize}
    \item \textbf{Количество функций:} Для $n$ переменных существует $2^{2^n}$ различных булевых функций.
    \item \textbf{Задание таблицы:} Функция задается таблицей истинности из $2^n$ строк. В учебнике Новикова принят \textbf{установленный порядок} перечисления наборов — лексикографический (по возрастанию двоичных чисел).
\end{itemize}

\subsection*{Существенные и фиктивные переменные}

\paragraph{Определение}
Переменная $x_i$ называется \textbf{существенной} для функции $f(x_1, \dots, x_n)$, если существует такой набор значений остальных переменных $(\alpha_1, \dots, \alpha_{i-1}, \alpha_{i+1}, \dots, \alpha_n)$, на котором изменение значения $x_i$ меняет значение функции:
$$f(\alpha_1, \dots, 0, \dots, \alpha_n) \neq f(\alpha_1, \dots, 1, \dots, \alpha_n).$$
В противном случае переменная называется \textbf{несущественной (фиктивной)}.

\paragraph{Удаление и введение переменных}
Булевы функции рассматриваются с точностью до несущественных переменных.
\begin{itemize}
    \item \textbf{Удаление:} Если $x_i$ фиктивна, таблицу можно сократить вдвое, удалив столбец $x_i$ и повторяющиеся строки.
    \item \textbf{Введение:} Любую функцию $f(x_1, \dots, x_n)$ можно рассматривать как функцию от большего числа переменных, добавив фиктивные.
\end{itemize}

\subsection*{Элементарные булевы функции}

\paragraph{Функции одной переменной ($n=1$)}
Существует $2^{2^1} = 4$ функции:
\begin{enumerate}
    \item \textbf{Константа 0} (Нуль).
    \item \textbf{Тождественная функция} $f(x) = x$.
    \item \textbf{Отрицание} $f(x) = \neg x$ (инверсия).
    \item \textbf{Константа 1} (Единица).
\end{enumerate}

\paragraph{Функции двух переменных ($n=2$)}
Существует $2^{2^2} = 16$ функций. Основные из них:
\begin{itemize}
    \item $x \land y$ — \textbf{Конъюнкция} (логическое умножение).
    \item $x \lor y$ — \textbf{Дизъюнкция} (логическое сложение).
    \item $x \oplus y$ — \textbf{Сложение по модулю 2} (XOR, исключающее ИЛИ).
    \item $x \to y$ — \textbf{Импликация} (если $x$, то $y$).
    \item $x \equiv y$ — \textbf{Эквивалентность}.
    \item $x | y$ — \textbf{Штрих Шеффера} (И-НЕ).
    \item $x \downarrow y$ — \textbf{Стрелка Пирса} (ИЛИ-НЕ).
\end{itemize}

\paragraph{Функции трех переменных ($n=3$)}
Существует $2^{2^3} = 256$ функций. Они обычно не имеют специальных названий и представляются в виде \textbf{формул} через функции 1 и 2 переменных.

